
\documentclass{article}
%%%%%%%%%%%%%%%%%%%%%%%%%%%%%%%%%%%%%%%%%%%%%%%%%%%%%%%%%%%%%%%%%%%%%%%%%%%%%%%%%%%%%%%%%%%%%%%%%%%%%%%%%%%%%%%%%%%%%%%%%%%%%%%%%%%%%%%%%%%%%%%%%%%%%%%%%%%%%%%%%%%%%%%%%%%%%%%%%%%%%%%%%%%%%%%%%%%%%%%%%%%%%%%%%%%%%%%%%%%%%%%%%%%%%%%%%%%%%%%%%%%%%%%%%%%%
\usepackage{amsfonts}

%TCIDATA{OutputFilter=LATEX.DLL}
%TCIDATA{Version=5.50.0.2890}
%TCIDATA{<META NAME="SaveForMode" CONTENT="1">}
%TCIDATA{BibliographyScheme=Manual}
%TCIDATA{Created=Saturday, July 29, 2006 18:07:07}
%TCIDATA{LastRevised=Friday, August 04, 2006 17:58:29}
%TCIDATA{<META NAME="GraphicsSave" CONTENT="32">}
%TCIDATA{<META NAME="DocumentShell" CONTENT="Standard LaTeX\Blank - Standard LaTeX Article">}
%TCIDATA{CSTFile=40 LaTeX article.cst}

\newtheorem{theorem}{Theorem}
\newtheorem{acknowledgement}[theorem]{Acknowledgement}
\newtheorem{algorithm}[theorem]{Algorithm}
\newtheorem{axiom}[theorem]{Axiom}
\newtheorem{case}[theorem]{Case}
\newtheorem{claim}[theorem]{Claim}
\newtheorem{conclusion}[theorem]{Conclusion}
\newtheorem{condition}[theorem]{Condition}
\newtheorem{conjecture}[theorem]{Conjecture}
\newtheorem{corollary}[theorem]{Corollary}
\newtheorem{criterion}[theorem]{Criterion}
\newtheorem{definition}[theorem]{Definition}
\newtheorem{example}[theorem]{Example}
\newtheorem{exercise}[theorem]{Exercise}
\newtheorem{lemma}[theorem]{Lemma}
\newtheorem{notation}[theorem]{Notation}
\newtheorem{problem}[theorem]{Problem}
\newtheorem{proposition}[theorem]{Proposition}
\newtheorem{remark}[theorem]{Remark}
\newtheorem{solution}[theorem]{Solution}
\newtheorem{summary}[theorem]{Summary}
\newenvironment{proof}[1][Proof]{\noindent\textbf{#1.} }{\ \rule{0.5em}{0.5em}}
\input{tcilatex}
\begin{document}


\subsection{\protect\bigskip Linear wave equation}

Electromagnetic fields are generated by charge and current densities
according to Maxwell's equations%
\begin{eqnarray}
\nabla \cdot \mathbf{E} &=&\frac{\rho }{\varepsilon _{0}}  \label{maxwell1}
\\
\nabla \cdot \mathbf{B} &=&0  \label{maxwell2} \\
\nabla \times \mathbf{E} &=&-\frac{\partial \mathbf{B}}{\partial t}
\label{maxwell3} \\
\nabla \times \mathbf{B} &=&\mu _{0}\mathbf{J}+\varepsilon _{0}\mu _{0}\frac{%
\partial \mathbf{B}}{\partial t}  \label{maxwell4}
\end{eqnarray}%
where $\mathbf{E}\left( \mathbf{X},t\right) $ and $\mathbf{B}\left( \mathbf{X%
},t\right) $ are the electric and magnetic fields, while $\rho \left( 
\mathbf{X},t\right) $ and $\mathbf{J}\left( \mathbf{X},t\right) $ are the
charge and current densities. Here, $\mathbf{X}$ is the vector position, and 
$t$ is the time. From the continuity equation, the charge density $\rho
\left( \mathbf{X},t\right) $ is linked to the current density $\mathbf{J}%
\left( \mathbf{X},t\right) $ by the relation 
\begin{equation}
\frac{\partial \rho }{\partial t}+\nabla \cdot \mathbf{J=0}
\label{continuity}
\end{equation}

The Maxwell equations \ref{maxwell1}-\ref{maxwell4} are a closed system only
if the charge and the current densities are known. In a material medium, $%
\rho \left( \mathbf{X},t\right) $ and $\mathbf{J}\left( \mathbf{X},t\right) $
are themselves functionals of $\mathbf{E}\left( \mathbf{X},t\right) $ and $%
\mathbf{B}\left( \mathbf{X},t\right) $. The expression of $\mathbf{J}\left( 
\mathbf{X},t\right) $ as a function of $\mathbf{E}\left( \mathbf{X},t\right) 
$ is named as the constitutive relation and may be expressed in general as a
convolution%
\begin{equation}
\mathbf{J}\left( \mathbf{X},t\right) =\diiint d^{3}\mathbf{X}^{\prime }\int
dt^{\prime }\mathbb{S}\left( \mathbf{X},\mathbf{X}^{\prime },t,t^{\prime
}\right) \cdot \mathbf{E}\left( \mathbf{X}^{\prime },t^{\prime }\right)
\label{constitutive}
\end{equation}%
where $\mathbb{S}\left( \mathbf{X},\mathbf{X}^{\prime },t,t^{\prime }\right) 
$ is the conductivity tensor which depends of the equilibrium properties of
the medium. It is assumed that electromagnetic fields have a small
amplitude, so that non-linear effects may be neglected. In a infinite plasma
with an homogeneous constant equilibrium which is invariant by translation
in space and time, the conductivity tensor is only a function of the
relative distance in space and time, and Eq. \ref{constitutive} may be
rewritten as%
\begin{equation}
\mathbf{J}\left( \mathbf{X},t\right) =\diiint d^{3}\mathbf{X}^{\prime }\int
dt^{\prime }\mathbb{S}\left( \mathbf{X}-\mathbf{X}^{\prime },t-t^{\prime
}\right) \cdot \mathbf{E}\left( \mathbf{X}^{\prime },t^{\prime }\right)
\label{constitutive2}
\end{equation}

In this case, fields and current can be Fourier analyzed, i.e. represented
as a superposition of time-harmonic plane waves%
\begin{equation}
\mathbf{E}\left( \mathbf{X},t\right) =\int d\omega \diiint d^{3}\mathbf{kE}_{%
\mathbf{k}}\left( \mathbf{k},\omega \right) e^{i\left( \mathbf{k\cdot X}%
-\omega t\right) }
\end{equation}%
\begin{equation}
\mathbf{J}\left( \mathbf{X},t\right) =\int d\omega \diiint d^{3}\mathbf{kJ}_{%
\mathbf{k}}\left( \mathbf{k},\omega \right) e^{i\left( \mathbf{k\cdot X}%
-\omega t\right) }
\end{equation}%
with the reverse relations%
\begin{equation}
\mathbf{E}_{\mathbf{k}}\left( \mathbf{k},\omega \right) =\diiint d^{3}%
\mathbf{X}\int dt\mathbf{E}\left( \mathbf{X},t\right) e^{-i\left( \mathbf{%
k\cdot X}-\omega t\right) }
\end{equation}%
\begin{equation}
\mathbf{J}_{\mathbf{k}}\left( \mathbf{k},\omega \right) =\diiint d^{3}%
\mathbf{X}\int dt\mathbf{J}\left( \mathbf{X},t\right) e^{-i\left( \mathbf{%
k\cdot X}-\omega t\right) }
\end{equation}%
and%
\begin{equation}
\mathbf{J}_{\mathbf{k}}\left( \mathbf{k},\omega \right) =\mathbb{S}\left( 
\mathbf{k},\omega \right) \cdot \mathbf{E}_{\mathbf{k}}\left( \mathbf{k}%
,\omega \right)  \label{constitutivefourier}
\end{equation}

Applying the Fourier transform to the Maxwell's equations lead to%
\begin{eqnarray}
i\mathbf{k}\cdot \mathbf{E}_{\mathbf{k}} &=&\frac{\rho _{\mathbf{k}}}{%
\varepsilon _{0}} \\
i\mathbf{k}\cdot \mathbf{B}_{\mathbf{k}} &=&0 \\
i\mathbf{k}\times \mathbf{E}_{\mathbf{k}} &=&i\omega \mathbf{B}_{\mathbf{k}}
\\
i\mathbf{k}\times \mathbf{B}_{\mathbf{k}} &=&\mu _{0}\mathbf{J}_{\mathbf{k}%
}-\varepsilon _{0}\mu _{0}i\omega \mathbf{B}_{\mathbf{k}}
\end{eqnarray}%
using the transformation $T.F\left( \nabla \right) =i\mathbf{k}.$

Combining the last two equations and using Eq. \ref{constitutivefourier}
leads to the linear wave equation%
\begin{equation}
\mathbf{k}\times \mathbf{k}\times \mathbf{E}_{\mathbf{k}}=-i\varepsilon
_{0}\mu _{0}\omega \mathbb{S}\left( \mathbf{k},\omega \right) \cdot \mathbf{E%
}_{\mathbf{k}}-\varepsilon _{0}\mu _{0}\omega ^{2}\mathbf{E}_{\mathbf{k}}
\label{Maxwellwaveeq}
\end{equation}%
or%
\begin{equation}
\mathbf{k}\times \mathbf{k}\times \mathbf{E}_{\mathbf{k}}+\frac{\omega ^{2}}{%
c^{2}}\mathbb{K}\left( \mathbf{k},\omega \right) \cdot \mathbf{E}_{\mathbf{k}%
}=0
\end{equation}%
where $\mathbb{K\left( \mathbf{k},\omega \right) }$ 
\begin{equation}
\mathbb{K\left( \mathbf{k},\omega \right) =I+X}\left( \mathbf{k},\omega
\right)  \label{permittivity}
\end{equation}
is the permittivity tensor, 
\begin{equation}
\mathbb{X\left( \mathbf{k},\omega \right) =}\frac{i\mathbb{S}\left( \mathbf{k%
},\omega \right) }{\varepsilon _{0}\mu _{0}}  \label{susceptibility}
\end{equation}
$\mathbb{S}\left( \mathbf{k},\omega \right) $ being the susceptibility
tensor and $\mathbb{I}$ the unitary matrix.

Introducing the index of refraction%
\begin{equation}
\mathbf{N\equiv }\frac{c}{\omega }\mathbf{k}  \label{indexrefraction}
\end{equation}%
the Maxwell equation becomes%
\begin{equation}
\mathbf{N}\times \mathbf{N}\times \mathbf{E}_{\mathbf{k}}+\mathbb{K}\left( 
\mathbf{k},\omega \right) \cdot \mathbf{E}_{\mathbf{k}}=0
\end{equation}%
and using the vectorial relation%
\begin{equation}
\mathbf{A}\times \mathbf{A}\times \mathbf{B=}\left( \mathbf{AA-A}^{2}\mathbb{%
I}\right) \cdot \mathbf{B}
\end{equation}%
where $\mathbf{AA}$ is the dyadic product and $\mathbf{A}^{2}=\mathbf{A}%
\cdot \mathbf{A}$, one obtains%
\begin{equation}
\left[ \mathbf{NN-N}^{2}\mathbb{I}+\mathbb{K}\left( \mathbf{k},\omega
\right) \right] \cdot \mathbf{E}_{\mathbf{k}}=0
\end{equation}%
which can be rewritten in the simple form%
\begin{equation}
\mathbb{M}\cdot \mathbf{E}_{\mathbf{k}}=0  \label{maxwellwaveeq1}
\end{equation}%
where $\mathbb{M}\left( \mathbf{k},\omega \right) $ 
\begin{equation}
\mathbb{M=}\mathbf{NN-N}^{2}\mathbb{I}+\mathbb{K}\left( \mathbf{k},\omega
\right) 
\end{equation}%
is the Maxwell or dispersion tensor, where $\mathbf{NN-N}^{2}\mathbb{I}$
part is called the vacuum term. Indeed, in the vacuum, $\mathbb{K=I.}$
Anticipating next section on the calculation of the susceptibility tensor $%
\mathbb{X}$, it can be shown that $\mathbb{X}$ depends from the following
non-dimensional parameters: the refractive index $\mathbf{N,}$ the thermal
velocity normalized to the speed of light $\beta _{th}=v_{th}/c,$ and for
each species $s$, the ratios $\overline{\omega }_{ps}=\omega _{ps}/\omega $
and $\overline{\Omega }=\Omega _{cs}/\omega $ of the plasma and cyclotron
frequencies $\omega _{ps}$ and $\Omega _{cs}$ to the wave frequency $\omega $%
. Therefore the Maxwell tensor may be expressed in the general form%
\begin{equation}
\mathbb{M=}\mathbf{NN-}\left( \mathbf{N}^{2}-1\right) \mathbb{I}+\mathbb{X}%
\left( \mathbf{N},\beta _{th},\overline{\omega }_{ps},\overline{\Omega }%
_{cs}\right) 
\end{equation}

In order to find a non-trival solution to the Eq. \ref{maxwellwaveeq1}, the
determinant of the Maxwell tensor must be zero, 
\begin{equation}
\mathcal{D\left( \mathbf{k},\omega \right) \equiv }\det \mathbb{M\left( 
\mathbf{k},\omega \right) =0}  \label{dispersionrelation}
\end{equation}%
which defines the dispersion relation. This relation gives the local
electromagnetic eigenmodes that cvan be excited independently in the
homogenous plasma. It can be solved for either $\omega \mathcal{\left( 
\mathbf{k}\right) }$ or $k_{i}\left( \omega ,k_{j},k_{k}\right) $ where $%
i,j,k$ is a permutation of coordinates of $\mathbf{k}$.

\subsection{Energy, power flux and ray equations for linear waves}

\bigskip 

\bigskip 

\bigskip

\subsection{Ray dynamics}

Ray evolution in anisotropic an inhomogeneous plasma is defined by the
condition  $\mathcal{D}\left( \mathbf{r},t,\mathbf{k,}\omega \right) \mathbb{%
=0}$ along the group velocity trajectory. Therefore, at all points along the
group velocity trajectory, the condition 
\begin{equation}
\delta \mathcal{D=}\frac{\partial \mathcal{D}}{\partial \mathbf{r}}\delta 
\mathbf{r}+\frac{\partial \mathcal{D}}{\partial t}\delta t\mathbf{+}\frac{%
\partial \mathcal{D}}{\partial \mathbf{k}}\delta \mathbf{k+}\frac{\partial 
\mathcal{D}}{\partial \omega }\delta \omega =0
\end{equation}%
must be satisfied. Let define $\tau $ a dimensionless parameter, which is a
measure of distance along the trajectory.  
\begin{equation}
\delta \mathcal{D=}\frac{\partial \mathcal{D}}{\partial \mathbf{r}}\mathbf{%
\dot{r}}\delta \tau +\frac{\partial \mathcal{D}}{\partial t}\dot{t}\delta
\tau \mathbf{+}\frac{\partial \mathcal{D}}{\partial \mathbf{k}}\mathbf{\dot{r%
}}\delta \tau \mathbf{+}\frac{\partial \mathcal{D}}{\partial \omega }\dot{%
\omega}\delta \tau =0  \label{deltaD2}
\end{equation}%
where $\mathbf{\dot{r}=}d\mathbf{r/}d\tau ,$ $\mathbf{\dot{k}=}d\mathbf{k/}%
d\tau ,$ $\dot{t}\mathbf{=}dt\mathbf{/}d\tau $ and $\dot{\omega}=d\omega 
\mathbf{/}d\tau .$ If the wave vector $\mathbf{k}$ may be expressed in
coordinates that are canonicaly conjugate to those of the position vector $%
\mathbf{r,}$ as well as the frequency $\omega $ for the time $t$,  the
relations%
\begin{eqnarray}
\frac{d\mathbf{r}}{d\tau } &=&\frac{\partial \mathcal{D}}{\partial \mathbf{k}%
} \\
\frac{d\mathbf{k}}{d\tau } &=&-\frac{\partial \mathcal{D}}{\partial \mathbf{r%
}} \\
\frac{dt}{d\tau } &=&-\frac{\partial \mathcal{D}}{\partial \omega } \\
\frac{d\omega }{d\tau } &=&\frac{\partial \mathcal{D}}{\partial t}
\end{eqnarray}
may be deduced from the Hamiltonian nature of $\mathcal{D}$, and Eq.\ref%
{deltaD2} is automatically satisfied. These four equations are therefore the
ray equations. From the first and the third equations, one deduce directly
the relation between the group velocity $\mathbf{v}_{g}$  
\begin{equation}
\frac{d\mathbf{r}}{dt}=\frac{\partial \mathcal{D}}{\partial \mathbf{k}}/%
\frac{\partial \mathcal{D}}{\partial \omega }=\frac{\partial \omega }{%
\partial \mathbf{k}}=\mathbf{v}_{g}
\end{equation}%
along with the wave energy flows. In the stationary case, i.e. $\partial 
\mathcal{D/}\partial t=0,$ $\omega $ is a constant of the ray dynamics, and
only the three first equations remain. If the time $t$ is considered as the
evolution parameter of the ray, two equations are necessary 
\begin{eqnarray}
\frac{d\mathbf{r}}{dt} &=&-\frac{\partial \mathcal{D}}{\partial \mathbf{k}}/%
\frac{\partial \mathcal{D}}{\partial \omega }  \label{rayeq1a} \\
\frac{d\mathbf{k}}{dt} &=&\frac{\partial \mathcal{D}}{\partial \mathbf{r}}/%
\frac{\partial \mathcal{D}}{\partial \omega }  \label{rayeq2a}
\end{eqnarray}%
which are the usual forms found in the litterature. However, one may also
keep the simple form%
\begin{eqnarray}
\frac{d\mathbf{r}}{d\tau } &=&\frac{\partial \mathcal{D}}{\partial \mathbf{k}%
}  \label{rayeq1b} \\
\frac{d\mathbf{k}}{d\tau } &=&-\frac{\partial \mathcal{D}}{\partial \mathbf{r%
}}  \label{rayeq2b}
\end{eqnarray}%
if the determination of the group velocity from $\partial \mathcal{D/}%
\partial \omega $ is not necessary. This reduces the number of derivative to
perform. 

For arbitrary tokamak magnetic equilibrium, the appropriate coordinate
system is $\left( \psi ,\theta ,\phi \right) $ for describing the position $%
\mathbf{r.}$ It is then necessary to find the canonically conjugate
coordinates of $\mathbf{k,}$ so that by definition%
\begin{equation}
\left\{ r_{i},k_{j}\right\} =\delta _{ij}
\end{equation}%
where $\left\{ \ldots \right\} $ is Poisson bracket, $\delta _{ij}$ is the
Kronecker delta, $r_{i}$ and $k_{j}$ being the coordinate of $\mathbf{r}$
and $\mathbf{k}$ respectively. For this purpose, all cross terms must be
zero, i.e.%
\begin{eqnarray}
\frac{\partial k_{\psi }}{\partial \psi } &=&\frac{\partial k_{\psi }}{%
\partial \theta }=\frac{\partial k_{\psi }}{\partial \phi }=0 \\
\frac{\partial k_{\theta }}{\partial \psi } &=&\frac{\partial k_{\theta }}{%
\partial \theta }=\frac{\partial k_{\theta }}{\partial \phi }=0 \\
\frac{\partial k_{\phi }}{\partial \psi } &=&\frac{\partial k_{\phi }}{%
\partial \theta }=\frac{\partial k_{\phi }}{\partial \phi }=0
\end{eqnarray}%
so all components of $\mathbf{k}$ must be independent of coordinates of $%
\mathbf{r.}$ Let defines canonically conjugate coordinates $\left( k_{\psi
},m,n\right) $ of $\mathbf{k}$, where $m$ and $n$ are the poloidal and
toroidal wave numbers respectively, the ray equations becomes 
\begin{eqnarray}
\frac{d\psi }{dt} &=&\frac{\partial \mathcal{D}}{\partial k_{\psi }}/\frac{%
\partial \mathcal{D}}{\partial \omega }  \label{h1} \\
\frac{d\theta }{dt} &=&\frac{\partial \mathcal{D}}{\partial m}/\frac{%
\partial \mathcal{D}}{\partial \omega }  \label{h2} \\
\frac{d\phi }{dt} &=&\frac{\partial \mathcal{D}}{\partial n}/\frac{\partial 
\mathcal{D}}{\partial \omega }  \label{h3}
\end{eqnarray}%
and 
\begin{eqnarray}
\frac{dk_{\psi }}{dt} &=&-\frac{\partial \mathcal{D}}{\partial \psi }/\frac{%
\partial \mathcal{D}}{\partial \omega }  \label{h4} \\
\frac{dm}{dt} &=&-\frac{\partial \mathcal{D}}{\partial \theta }/\frac{%
\partial \mathcal{D}}{\partial \omega }  \label{h5} \\
\frac{dn}{dt} &=&-\frac{\partial \mathcal{D}}{\partial \phi }/\frac{\partial 
\mathcal{D}}{\partial \omega }  \label{h6}
\end{eqnarray}

For an axisymmetric magnetic configuration, $n$ is an invariant, so that
from Eq.\ref{h6}, $dn/dt=0$. Therefore, in this case, $\partial \mathcal{D}%
/\partial \phi =0,$ and only five equations must be solved. The value of $n$
is given by conditions at launch, i.e. antenna properties.

When the determination of the time evolution of $\mathbf{r}$ and $\mathbf{k}$
is useless, the time variable $t$ must be replaced by $\tau ,$ and $\partial
D/\partial \omega =1$ must be enforced in the relations \ref{h1}-\ref{h6}.

Though it is possible to keep the general coordinate systems for the
dispersion relation, the case of a plasma with a uniform static magnetic
field led to define usualy the parallel $k_{\Vert }$ and the perpendical $%
k_{\bot }$ components of $\mathbf{k}$ with respect to the magnetic field
direction $\widehat{\mathbf{b}}=\mathbf{B}/\left\Vert \mathbf{B}\right\Vert .
$

For determining $\mathbf{k}^{2}$ and the scalar product $k_{\Vert }=\mathbf{%
k\cdot }\widehat{\mathbf{b}}$ in the general coordinates system, it is
necessary to introduce the general scalar product $\mathbf{x}\cdot \mathbf{y}%
=g^{ij}x_{i}y_{j}$ where $g^{ij}$ is the metric tensor 
\begin{equation}
g^{ij}=\left( 
\begin{array}{ccc}
\left\vert \nabla \psi \right\vert ^{2} & -\frac{\left\vert \nabla \psi
\right\vert \sin \alpha }{r} & 0 \\ 
-\frac{\left\vert \nabla \psi \right\vert \sin \alpha }{r} & 1/r^{2} & 0 \\ 
0 & 0 & 1/R^{2}%
\end{array}%
\right) 
\end{equation}%
with the covariant components $x_{i}$ of the vector $\mathbf{x}$%
\begin{equation}
\mathbf{x}=x_{i}\mathbf{e}^{i}
\end{equation}%
$\mathbf{e}^{i}$ being the contravariant vector basis. Here, $\cos \alpha =%
\widehat{\mathbf{\psi }}\cdot \widehat{\mathbf{r}}.$ 

Since the magnetic field $\mathbf{B}$ is tangent to a poloidal flux surface $%
\psi $, it is defined by its poloidal $B_{p}$ and toroidal $B_{T}$
components only. In the general coordinate system, 
\begin{equation}
\mathbf{B=}B_{p}\widehat{\mathbf{s}}+B_{T}\widehat{\mathbf{\phi }}
\end{equation}%
where $\widehat{\mathbf{s}}$ is tangent to the flux surface and $%
B=\left\Vert \mathbf{B}\right\Vert =\sqrt{B_{T}^{2}+B_{p}^{2}},$ so that  
\begin{equation}
\widehat{\mathbf{b}}\mathbf{=}\cos \alpha \frac{B_{p}}{B}\widehat{\mathbf{%
\theta }}+\frac{B_{T}}{B}\widehat{\mathbf{\phi }}
\end{equation}%
or%
\begin{eqnarray}
\widehat{\mathbf{b}} &\mathbf{=}&r\cos \alpha \frac{B_{p}}{B}\left( \frac{%
\widehat{\mathbf{\theta }}}{r}\right) +R\frac{B_{T}}{B}\left( \frac{\widehat{%
\mathbf{\phi }}}{R}\right)   \nonumber \\
&=&r\cos \alpha \frac{B_{p}}{B}\mathbf{e}^{\theta }+R\frac{B_{T}}{B}\mathbf{e%
}^{\phi }
\end{eqnarray}%
where the contravariant vector basis is $\left( \mathbf{e}^{\psi },\mathbf{e}%
^{\theta },\mathbf{e}^{\phi }\right) =\left( \left\vert \nabla \psi
\right\vert \widehat{\mathbf{\psi }},\widehat{\mathbf{\theta }}/r,\widehat{%
\mathbf{\phi }}/R\right) .$ 

Expressions \ref{h4}-\ref{h6}, may be expressed in a straighforward manner
in a vectorial form%
\begin{equation}
\frac{d\mathbf{k}}{dt}=\nabla \mathcal{D}/\frac{\partial \mathcal{D}}{%
\partial \omega }
\end{equation}%
and using the gradient definition for the general system of coordinate here
considered,%
\begin{eqnarray}
\frac{d\mathbf{k}}{dt} &=&\frac{\partial \mathcal{D}}{\partial \psi }/\frac{%
\partial \mathcal{D}}{\partial \omega }\mathbf{e}^{\psi }+\frac{\partial 
\mathcal{D}}{\partial \theta }/\frac{\partial \mathcal{D}}{\partial \omega }%
\mathbf{e}^{\theta }+\frac{\partial \mathcal{D}}{\partial \phi }/\frac{%
\partial \mathcal{D}}{\partial \omega }\mathbf{e}^{\phi } \\
&=&\frac{dk_{\psi }}{dt}\mathbf{e}^{\psi }+\frac{dm}{dt}\mathbf{e}^{\theta }+%
\frac{dn}{dt}\mathbf{e}^{\phi }  \nonumber
\end{eqnarray}%
so that canonically conjugate covariant coordinates of $\mathbf{k}$ are
simply $\left( k_{\psi },m,n\right) .$ From the coordinates of $\widehat{%
\mathbf{b}}$ and $\mathbf{k,}$ 

\begin{eqnarray}
k_{\Vert } &=&\frac{\mathbf{k\cdot B}}{\left\vert \mathbf{B}\right\vert }=%
\mathbf{k\cdot }\widehat{\mathbf{b}}  \nonumber \\
&=&\left( -\left\vert \nabla \psi \right\vert k_{\psi }\sin \alpha +\frac{m}{%
r}\right) \cos \alpha \frac{B_{p}}{B}+\frac{n}{R}\frac{B_{T}}{B}
\label{kpar}
\end{eqnarray}%
and     
\begin{equation}
\mathbf{k}^{2}=\left\vert \nabla \psi \right\vert ^{2}k_{\psi }^{2}-2\frac{%
\left\vert \nabla \psi \right\vert \sin \alpha }{r}mk_{\psi }+\frac{m^{2}}{%
r^{2}}+\frac{n^{2}}{R^{2}}  \label{k2}
\end{equation}%
with $\left\vert \nabla \psi \right\vert =RB_{P}.$ It is interesting to
notice that%
\begin{equation}
\nabla \psi \cdot \nabla \theta =\frac{\left\vert \nabla \psi \right\vert }{r%
}\left( \widehat{\psi }\cdot \widehat{\theta }\right) =-\frac{\left\vert
\nabla \psi \right\vert }{r}\sin \alpha 
\end{equation}%
from geometrical definitions of the angle $\theta .$

For a circular plasma cross-section, $\cos \alpha =\widehat{\psi }\cdot 
\widehat{r}=1,$ and $\sin \alpha =0.$ In this limit, the intercoupling in $%
\mathbf{k}^{2}$ between $m$ and $k_{\psi }$ which results from the poloidal
plasma shaping vanishes$.$

From Eq.\ref{kpar}, the usual expression for circular plasma cross-section%
\begin{equation}
k_{\Vert }^{\func{Ci}rc.}=+\frac{m}{r}\frac{B_{p}}{B}+\frac{nB_{T}}{RB}
\end{equation}%
is recovered, and using the general expression of the local safety factor $%
q\left( \psi ,\theta \right) $%
\begin{equation}
q\left( \psi ,\theta \right) =\frac{I\left( \psi \right) J}{R^{2}}
\end{equation}%
where the poloidal current $I\left( \psi \right) =RB_{T}$ and the inverse of
the Jacobian is $J^{-1}=$ $\left( \widehat{\psi }\cdot \widehat{r}\right)
B_{P}/r,$ with $\widehat{\psi }\cdot \widehat{r}=1,$

\begin{eqnarray}
k_{\Vert }^{\func{Ci}rc.} &=&\frac{1}{B}\left( m\frac{B_{p}}{r}+n\frac{B_{T}%
}{R}\right)   \nonumber \\
&=&\frac{q}{JB}\left( \frac{m}{q}+n\right)   \nonumber \\
&=&\frac{1}{R}\frac{B_{T}}{B}\left( \frac{m}{q}+n\right) 
\end{eqnarray}

In the small inverse aspect ratio limit $\epsilon =r/R_{p}\ll 1$, $R\simeq
R_{p}$ and $B=B_{T},$ so that the traditional form of $k_{\Vert }$

\begin{equation}
k_{\Vert }^{\func{Ci}rc.}\left( \epsilon \ll 1\right) \simeq \frac{1}{R_{p}}%
\left( \frac{m}{q}+n\right) 
\end{equation}%
is well recovered.

\subsection{Ray equations}

Regarding the complexity of the expressions of $k_{\Vert }$ and $k_{\bot },$
it is more efficient numerically to express directly the dispersion relation
in the general system of coordinates for $\mathbf{r}$ and $\mathbf{k.}$Using
the relation $\mathbf{N=}\left( \omega /c\right) \mathbf{k}$,  

\begin{equation}
\mathbb{M=}\frac{c^{2}}{\omega ^{2}}\mathbf{kk-}\left( \frac{c^{2}}{\omega
^{2}}\mathbf{k}^{2}-1\right) \mathbb{I}+\mathbb{X}\left( \mathbf{N},\beta
_{th},\overline{\omega }_{ps},\overline{\Omega }_{cs}\right) 
\end{equation}%
where $\mathbf{k}^{2}$ is given by relation \ref{k2}.

Using the metric, 

\bigskip

\bigskip

\bigskip 

\bigskip 

\bigskip 

\bigskip 

In a uniform plasma where an external static magnetic field is applied, 
\begin{equation}
\mathbf{B=}\overline{B}\widehat{\mathbf{b}}
\end{equation}%
where $\overline{B}=\left\Vert \mathbf{B}\right\Vert ,$ the cylindrical
symmetry around the magnetic field direction $\widehat{\mathbf{b}}$ imposes
that the dispersion relation be only a function of the parallel and
perpendicular components of the wave vector $\mathbf{k}$%
\begin{eqnarray}
k_{\Vert } &=&\mathbf{k\cdot }\widehat{\mathbf{b}} \\
k_{\bot } &=&\left\Vert \mathbf{k-k\cdot }\widehat{\mathbf{b}}\right\Vert
\end{eqnarray}

In addition, since all directions perpendicular to the static magnetic field
are equivalent, the dispersion relation can be then written as $\mathcal{%
D\left( k_{\Vert },k_{\bot }^{2},\omega \right) =0.}$

The group velocity $\mathbf{v}_{g}$ defined as%
\begin{equation}
\mathbf{v}_{g}=\left. \frac{\partial \omega }{\partial \mathbf{k}}%
\right\vert _{\mathcal{D\left( \mathbf{k},\omega \right) }\mathbb{=0}}
\label{groupvelocity}
\end{equation}%
must lie in the $\left( \mathbf{k,}\widehat{\mathbf{b}}\right) $ plane. Note
that the phase velocity $\mathbf{v}_{p}$%
\begin{equation}
\mathbf{v}_{p}=\left. \frac{\omega }{k}\widehat{\mathbf{k}}\right\vert _{%
\mathcal{D\left( \mathbf{k},\omega \right) }\mathbb{=0}}
\label{phasevelocity}
\end{equation}%
where $\widehat{\mathbf{k}}=\mathbf{k}/k$ and $k=\left\Vert \mathbf{k}%
\right\Vert ,$ is not necessarily equal to the group velocity, since in
dispersive medium, $\omega $ and $k$ are not proportional in general. In
addition, in anisotropic medium, i.e. when a static magnetic field is
applied, $\mathbf{v}_{g}$ is usualy not colinear to the wave vector
direction $\widehat{\mathbf{k}}.$

\bigskip

\bigskip

\subsection{Hot plasma dielectric tensor}

Using the relation $\mathbb{K=I+X}_{M}\mathbb{,}$

\begin{equation}
\mathbb{D_{M}=}\mathbf{NN}-\left( N^{2}-1\right) \mathbb{I}+\mathbb{X}_{M}
\end{equation}%
where $\ \mathbb{X}_{M}$ is the hot plasma susceptibility tensor for a
non-relativistic plasma which a the sum over contributions of all species 
\begin{equation}
\mathbb{X}_{M}\left( \mathbf{k},\omega \right) =\sum_{s}\mathbb{X}%
_{M}^{s}\left( \mathbf{k},\omega \right)
\end{equation}%
$\mathbb{X}_{M}^{s}\left( \mathbf{k},\omega \right) $ being itself the sum
over contributions from all harmonics numbers $n$%
\begin{equation}
\mathbb{X}_{M}^{s}\left( \mathbf{k},\omega \right) =\frac{\omega _{ps}^{2}}{%
\omega ^{2}}\zeta _{0s}\sum_{n=-\infty }^{n=+\infty }\mathbb{Y}%
_{M}^{s,n}\left( \mathbf{k},\omega \right)
\end{equation}

Matrix coefficients of $\mathbb{Y}_{M}^{s,n}$ are%
\begin{eqnarray}
\mathbb{Y}_{M,xx}^{s,n} &=&\frac{n^{2}}{\lambda _{s}}\Gamma _{ns}Z_{ns}
\label{a} \\
\mathbb{Y}_{M,yy}^{s,n} &=&\left( \frac{n^{2}}{\lambda _{s}}\Gamma
_{ns}-2\lambda _{s}n^{2}\Gamma _{ns}^{\prime }\right) Z_{ns}  \label{b} \\
\mathbb{Y}_{M,zz}^{s,n} &=&-\Gamma _{ns}\zeta _{ns}Z_{ns}^{\prime }
\label{c} \\
\mathbb{Y}_{M,xy}^{s,n} &=&i\sigma _{s}n\Gamma _{ns}^{\prime }Z_{ns}
\label{d} \\
\mathbb{Y}_{M,yx}^{s,n} &=&-\mathbb{Y}_{M,xy}^{s,n}  \label{e} \\
\mathbb{Y}_{M,xz}^{s,n} &=&-\sigma _{\Vert }\frac{n}{\sqrt{2\lambda _{s}}}%
\Gamma _{ns}Z_{ns}^{\prime }  \label{f} \\
\mathbb{Y}_{M,xz}^{s,n} &=&\mathbb{Y}_{M,zx}^{s,n}  \label{g} \\
\mathbb{Y}_{M,yz}^{s,n} &=&i\sigma _{s}\sigma _{\Vert }\sqrt{\frac{\lambda
_{s}}{2}}\Gamma _{ns}^{\prime }Z_{ns}^{\prime }  \label{h} \\
\mathbb{Y}_{M,zy}^{s,n} &=&-\mathbb{Y}_{M,yz}^{s,n}  \label{i}
\end{eqnarray}%
with%
\begin{eqnarray*}
\Gamma _{ns} &=&e^{-\lambda _{s}}I_{n}\left( \lambda _{s}\right) \\
Z_{ns} &=&Z\left( \zeta _{ns}\right) \\
\lambda _{s} &=&\left( \frac{k_{\bot }v_{Ts}}{\omega _{cs}}\right)
^{2}=\left( k_{\bot }\rho _{Ts}\right) ^{2} \\
\zeta _{ns} &=&\frac{\omega -\omega _{cs}}{\left\vert k_{\Vert }\right\vert
v_{Ts}} \\
\sigma _{s} &=&sign\left( q_{s}\right) \\
\sigma _{\Vert } &=&sign\left( k_{\Vert }\right)
\end{eqnarray*}%
where $I_{n}\left( x\right) $ is the modified Bessel function of the first
kind, $Z\left( x\right) $ is the plasma dispersion function, $\rho _{Ts}$ is
the thermal Larmor radius for the species $s$.

The plasma dispersion function $Z\left( x\right) $ is given by the relation%
\begin{equation}
Z\left( x\right) =i\sqrt{\pi }Erfc\left( ix\right)
\end{equation}%
for $\func{Im}\left( x\right) >0,$ or in power series in the form%
\begin{equation}
Z\left( x\right) =e^{-x^{2}}\left[ i\sqrt{\pi }-x\sum_{n=0}^{\infty }\frac{%
x^{2n}}{n!\left( 2n+1\right) }\right]
\end{equation}%
Ray dynamics in arbitrary magnetic equilibrium

\end{document}
